\RequirePackage[l2tabu, orthodox]{nag}

\documentclass[german,a4paper,fontsize=10pt,version=last]{scrlttr2}%german,

\input{../header.tex}
%\selectlanguage{german}
\DTMsetup{useregional=text}%numeric
\LoadLetterOption{DIN}% reload to make consistent again after loading geometry 
% else text does not reach lower boundary 
\input{descSender.tex}

\newcommand{\subject}{Meilenstein Dokumentation }
\input{headerLetter.tex}


\input{../headerSuppressMetaPDF.tex}

\hypersetup{
  pdfinfo={
    Author      ={Ernst Reissner},
    Title       ={Letter with/about the latex-maven-plugin },
    Subject     ={Brief von und ueber scrlttr2},
    Keywords    ={LaTeX;scrlttr2}
  }
}


\begin{document}

\newcommand{\AdressatFirma}{CapMini GmbH}
\newcommand{\AdressatTitel}{}
\newcommand{\AdressatVorname}{Egon}
\newcommand{\AdressatNachname}{Walther}
\newcommand{\AdressatAnschriftStrasse}{Walthershofener Straße 20}
\newcommand{\AdressatAnschriftPLZOrt}{88444 Walthershofen}
\setboolean{isMale}{true}

\begin{letter}{\Anschrift} 
  %\opening{Lieber Egon,}
  \opening{\Anrede}
%'\localeinfo{language.tag.bcp47}'
  Wie Sie wissen, haben wir einen Teil des Projektantrages in \LaTeX{} verfasst. 
  Nun habe ich alle Dokumentation bei jedem \foreignlanguage{english}{Build} 
  der Software oder der zugehörigen Seite 
  durch ein \foreignlanguage{english}{Maven Plugin} aus \LaTeX{} generiert. 
  Leider basiert das \foreignlanguage{english}{Plugin} auf sehr alten Techniken, 
  sodass ich zwar meine Diplomarbeit kompilieren kann, 
  nicht aber die \LaTeX{} Dokumente des Projektantrages. 

  Die Version 2.0 des Plugins bringt uns dem Ziel näher, 
  alle Dokumente, auch die nach neuester Technik zu kompilieren. 
  Zudem ist es jetzt auch möglich in qualitativ akzeptabler Weise 
  Briefe zu erstellen. 
  Insbesondere dieser Brief ist mit dem Plugin generiert. 

  %\lipsum{1}
  


  \closing{\Grussworte\\[2\baselineskip]}

\end{letter}

\end{document}